% ================= Topic 3 (EN) -- part d: problems 41--50 (challenge) =================
\begin{probbox}{41}{\stars{1}}
On a square dot grid, Ruby stands at dot $A$. She tosses a fair coin $6$ times; each time it shows heads she moves \emph{up} one dot, and
each time it shows tails she moves \emph{right} one dot (as in problem 16). The dots $M$, $N$, $B$ are marked in the figure.
\begin{figC}\begin{scope}[scale=0.85]
  \foreach \x in {0,...,3}{\foreach \y in {0,...,3}{\fill[Navy] (\x,\y) circle (2.4pt);}}
  \fill[tmOrange] (0,0) circle (3.4pt);\node[below left] at (0,0) {$A$};
  \fill[tmGreen!80!black] (1,1) circle (3.4pt);\node[above left] at (1,1) {$M$};
  \fill[AnsRed] (2,1) circle (3.4pt);\node[below right] at (2,1) {$N$};
  \fill[tmPurple] (3,3) circle (3.4pt);\node[above right] at (3,3) {$B$};
\end{scope}\end{figC}
\begin{enumerate}[label=\alph*),leftmargin=1.6em,itemsep=1pt,topsep=2pt]
\item Find the probability that after $6$ tosses Ruby stops at dot $B$.
\item Find the probability that Ruby passes through dot $M$ \emph{and} stops at $B$.
\item Find the probability that Ruby stops at $B$ \emph{without} passing through dot $N$.
\end{enumerate}
\end{probbox}
\needspace{8\baselineskip}
\probsol{
\textbf{Understanding the problem.} Each sequence of $6$ tosses is a sequence of $6$ symbols $\uparrow$ (up) or $\rightarrow$ (right).
There are $2^6=2\times2\times2\times2\times2\times2=64$ sequences, all equally likely. Each sequence is a \emph{path} on the grid.
We label positions by the pair (number of moves right, number of moves up): $A=(0,0)$, $M=(1,1)$, $N=(2,1)$, $B=(3,3)$.

\step{1}{Count the paths from $A$ to every dot (part a)}
\begin{recallbox}A dot can only be reached from the dot just to its left or the dot just below it. Therefore:
number of paths to a dot $=$ number of paths to the dot on its left $+$ number of paths to the dot below (addition principle).\end{recallbox}
Write the number of paths at each dot, starting from $A$ (write $1$); the bottom row and the left column are all $1$:
\begin{figC}\begin{scope}[scale=0.85]
  \foreach \x/\y/\n in {0/0/1,1/0/1,2/0/1,3/0/1,0/1/1,1/1/2,2/1/3,3/1/4,0/2/1,1/2/3,2/2/6,3/2/10,0/3/1,1/3/4,2/3/10,3/3/20}{
    \fill[tmBlue!12] (\x,\y) circle (9pt);\node[font=\small\bfseries,text=Navy] at (\x,\y) {\n};}
  \draw[tmPurple,thick] (3,3) circle (10pt);\node[right=9pt,text=tmPurple] at (3,3) {$B$};
\end{scope}\end{figC}
There are $20$ paths from $A$ to $B$. (Another way: reaching $B$ needs exactly $3$ up-moves among the $6$ moves; the number of ways to
choose $3$ turns out of $6$ is $\dfrac{6\times5\times4}{3\times2\times1}=20$.)
\[P(a)=\frac{20}{64}=\frac{5}{16}.\]

\step{2}{Part b: through $M$ and then to $B$ -- multiplication principle}
A path through $M$ has two stages: stage $A\to M$ ($1$ up, $1$ right) and stage $M\to B$ ($2$ up, $2$ right).
Stage $A\to M$ can be done in $2$ ways ($\uparrow\rightarrow$ or $\rightarrow\uparrow$); stage $M\to B$ in $6$ ways (as in problem 16:
choose $2$ up-turns out of $4$). Each way of doing the first stage combined with each way of doing the second gives a different path:
$2\times6=12$ paths.
\[P(b)=\frac{12}{64}=\frac{3}{16}.\]

\step{3}{Part c: avoiding $N$ -- use the complement}
Paths from $A$ through $N$ to $B$: stage $A\to N$ ($2$ right, $1$ up) has $3$ ways; stage $N\to B$ ($1$ right, $2$ up) has $3$ ways;
in total $3\times3=9$ paths. So the number of paths from $A$ to $B$ \emph{not} through $N$ is $20-9=11$.
\[P(c)=\frac{11}{64}.\]
\begin{proofbox}\emph{Check by ``blocking'' dot $N$.} Recompute the path numbers but put $0$ at $N$ (as if it cannot be used).
Row $1$: $1,\ 2,\ \mathbf{0},\ 0+1=1$. Row $2$: $1,\ 3,\ 3+0=3,\ 3+1=4$. Row $3$: $1,\ 4,\ 4+3=7,\ 7+4=11$.
At $B$ we get $11$, matching the complement.\end{proofbox}
\begin{warnbox}Parts b and c ask for probabilities among \emph{all} $64$ toss sequences, not among the $20$ paths to $B$. Dividing by $20$
answers a different question (a conditional probability, not yet studied).\end{warnbox}
\verify{In part a the possible final dots are $(6,0),(5,1),\dots,(0,6)$ with $1,6,15,20,15,6,1$ paths; the total is $64$. Dot $B=(3,3)$
corresponds to $20$.}
\ans{a) $\dfrac5{16}$;\quad b) $\dfrac3{16}$;\quad c) $\dfrac{11}{64}$}
}

\begin{probbox}{42}{\stars{2}}
We build grids of squares from toothpicks as in problem 17. An ``$m\times n$'' grid has $m$ rows, each with $n$ squares.
\begin{enumerate}[label=\alph*),leftmargin=1.6em,itemsep=1pt,topsep=2pt]
\item Find a formula for the number of toothpicks needed for an $m\times n$ grid.
\item A square $k\times k$ grid uses exactly $760$ toothpicks. Find $k$.
\item Sam has exactly $2026$ toothpicks and wants to build an $m\times n$ grid (with $1\le m\le n$) that uses \emph{all} of them.
Find all the grids Sam can build.
\end{enumerate}
\end{probbox}
\probsol{
\textbf{Understanding the problem.} As in problem 17, we count horizontal and vertical toothpicks separately. Part c leads to an equation in
two natural-number unknowns; the trick ``multiply by $2$ and add $1$'' turns it into a problem about divisors.

\step{1}{Part a: the formula}
An $m\times n$ grid has $m+1$ horizontal lines of $n$ toothpicks each: $(m+1)\times n$ horizontal toothpicks. It has $n+1$ vertical lines of
$m$ toothpicks each: $(n+1)\times m$ vertical toothpicks. Total:
\[T=(m+1)n+(n+1)m=mn+n+mn+m=2mn+m+n.\]
Check with the $1\times10$ grid: $2\times1\times10+1+10=31$, as in problem 17.

\step{2}{Part b: a square $k\times k$ grid}
Put $m=n=k$: $T=2k^2+2k=2k(k+1)$. We need $2k(k+1)=760$, that is $k(k+1)=380$.
$k(k+1)$ is the product of two consecutive natural numbers. Since $19\times20=380$, $k=19$. (As $k$ grows, $k(k+1)$ grows, so there is only
one value: $18\times19=342<380<420=20\times21$.)

\step{3}{Part c: set up the equation}
We need natural numbers $1\le m\le n$ with $2mn+m+n=2026$.
\begin{tipbox}Multiply both sides by $2$ and add $1$: $4mn+2m+2n+1=4053$. The left side factorises:
$(2m+1)(2n+1)=4mn+2m+2n+1$ (expand to check). So
\[(2m+1)(2n+1)=4053.\]\end{tipbox}

\step{4}{Prime factorisation of $4053$}
The digit sum of $4053$ is $4+0+5+3=12$, divisible by $3$: $4053=3\times1351$. Next, $1351=7\times193$.
The number $193$ is prime: it is not divisible by $2,3,5,7,11,13$, and $13\times13=169<193<196=14\times14$, so no more trials are needed.
\begin{recallbox}To check whether $N$ is prime, it is enough to try dividing by the primes $p$ with $p\times p\le N$.\end{recallbox}
So $4053=3\times7\times193$. The divisors of $4053$ are: $1,\ 3,\ 7,\ 21,\ 193,\ 579,\ 1351,\ 4053$.

\step{5}{Find the pairs $(m,n)$}
Since $m\ge1$, $2m+1\ge3$; since $m\le n$, $2m+1\le2n+1$. The ways of writing $4053$ as a product of two such divisors:
\begin{center}\begin{tabular}{|c|c|c|c|c|}\hline
$2m+1$ & $2n+1$ & $m$ & $n$ & Check $2mn+m+n$\\\hline
$3$ & $1351$ & $1$ & $675$ & $1350+1+675=2026$\\\hline
$7$ & $579$ & $3$ & $289$ & $1734+3+289=2026$\\\hline
$21$ & $193$ & $10$ & $96$ & $1920+10+96=2026$\\\hline
\end{tabular}\end{center}
(The pair $1\times4053$ is rejected because it gives $m=0$.)
\begin{warnbox}Trying $m=1,2,3,\dots$ one by one also works but is very long and it is easy to miss a solution. The ``multiply by $2$, add $1$''
trick turns the problem into ``find the divisors'' and guarantees that \emph{no} solution is missed.\end{warnbox}
\verify{For any other value of $m$, $2m+1$ is not a divisor of $4053$, so $n$ is not a natural number. For example $m=2$: $5(2n+1)=4053$ is
impossible because $4053$ is not divisible by $5$.}
\ans{a) $2mn+m+n$;\quad b) $k=19$;\quad c) three grids: $1\times675$, $3\times289$, $10\times96$}
}

\begin{probbox}{43}{\stars{1}}
Three fair dice coloured red, blue and yellow are rolled together (each is numbered $1$ to $6$). Find the probability of each event:
\begin{enumerate}[label=\alph*),leftmargin=1.6em,itemsep=1pt,topsep=2pt]
\item the three numbers are all different;
\item the number on the red die is less than the number on the blue die, and the number on the blue die is less than the number on the
yellow die;
\item at least one die shows a $6$;
\item the largest of the three numbers is exactly $4$.
\end{enumerate}
\end{probbox}
\probsol{
\textbf{Understanding the problem.} Each outcome is a triple (red, blue, yellow). By the multiplication principle there are
$6\times6\times6=216$ equally likely outcomes.

\step{1}{Part a: three different numbers}
Red die: $6$ ways; blue die different from red: $5$ ways; yellow die different from both: $4$ ways. That gives $6\times5\times4=120$ outcomes.
\[P(a)=\frac{120}{216}=\frac{5}{9}.\]

\step{2}{Part b: red $<$ blue $<$ yellow -- use symmetry}
Take a set of three different numbers, for example $\{2,3,5\}$. There are $3\times2\times1=6$ ways to assign these three numbers to the three
dice ($235$, $253$, $325$, $352$, $523$, $532$), and \emph{exactly one} of them is increasing ($2<3<5$).
So the $120$ outcomes of part a split into groups of $6$, each with $1$ favourable outcome: there are $120\div6=20$ favourable outcomes.
\[P(b)=\frac{20}{216}=\frac{5}{54}.\]
\begin{recallbox}The number of ways to arrange $3$ different objects in $3$ places is $3\times2\times1=6$.\end{recallbox}

\step{3}{Part c: at least one $6$ -- complement}
The opposite event is ``no die shows $6$'': each die has only $5$ outcomes ($1$ to $5$), in total $5\times5\times5=125$.
So the number of outcomes with at least one $6$ is $216-125=91$.
\[P(c)=\frac{91}{216}.\]
\begin{warnbox}Do not argue ``each die shows $6$ with probability $\frac16$, so for three dice it is $\frac36=\frac12$''. That addition counts
twice the outcomes with two or three sixes (for example $(6,6,1)$ is counted twice).\end{warnbox}

\step{4}{Part d: the largest number is exactly $4$}
``Largest number $\le4$'' means all three numbers are in $\{1,2,3,4\}$: $4\times4\times4=64$ outcomes.
``Largest number $\le3$'': $3\times3\times3=27$ outcomes. The largest number is exactly $4$ when it is $\le4$ but not $\le3$:
$64-27=37$ outcomes.
\[P(d)=\frac{37}{216}.\]
\begin{proofbox}Every outcome whose largest number is $\le3$ also belongs to the group whose largest number is $\le4$. Removing the smaller
group leaves exactly the outcomes whose largest number is at most $4$ but more than $3$, that is equal to $4$.\end{proofbox}
\verify{The idea of part d also gives part c: largest number equal to $6$ in $6^3-5^3=216-125=91$ outcomes, a match. The numbers of
outcomes by largest number $1,2,\dots,6$: $1+7+19+37+61+91=216$.}
\ans{a) $\dfrac59$;\quad b) $\dfrac5{54}$;\quad c) $\dfrac{91}{216}$;\quad d) $\dfrac{37}{216}$}
}

\begin{probbox}{44}{\stars{2}}
Special six-sided dice are made with one of the numbers $1,2,3,4,5,6$ written on each face (repeats allowed), so that, as in
problem 19, the probability of rolling a multiple of $3$ is $\frac12$ and the probability of rolling an even number is $\frac13$.
Two dice are considered the same if they carry the same six numbers (regardless of the order of the faces).
\begin{enumerate}[label=\alph*),leftmargin=1.6em,itemsep=1pt,topsep=2pt]
\item What is the largest possible sum of the six numbers on such a die? What is the smallest?
\item How many different such dice are there?
\end{enumerate}
\end{probbox}
\probsol{
\textbf{Understanding the problem.} As in problem 19, the conditions mean: \textbf{exactly $3$ faces are multiples of $3$} and
\textbf{exactly $2$ faces are even}. The key point is that $6$ belongs to both groups, so we sort the numbers $1$--$6$ into four types:
\begin{center}\begin{tabular}{|c|c|c|}\hline
Type & Numbers & Property\\\hline
I & $6$ & even and a multiple of $3$\\\hline
II & $3$ & odd, multiple of $3$\\\hline
III & $2,\ 4$ & even, not a multiple of $3$\\\hline
IV & $1,\ 5$ & odd, not a multiple of $3$\\\hline
\end{tabular}\end{center}

\step{1}{Express the number of faces of each type in terms of the number of $6$s}
Let $a$ be the number of faces showing $6$. Then:
\begin{itemize}[leftmargin=1.2em,itemsep=1pt,topsep=2pt]
\item the number of multiples of $3$ is (type I) $+$ (type II) $=3$, so type II has $3-a$ faces;
\item the number of even faces is (type I) $+$ (type III) $=2$, so type III has $2-a$ faces;
\item type IV has $6-a-(3-a)-(2-a)=1+a$ faces.
\end{itemize}
A number of faces cannot be negative, so $2-a\ge0$, that is $a\in\{0,1,2\}$.
\begin{center}\begin{tabular}{|c|c|c|c|c|}\hline
$a$ & type I (the $6$s) & type II (the $3$s) & type III ($2$ or $4$) & type IV ($1$ or $5$)\\\hline
$0$ & $0$ & $3$ & $2$ & $1$\\\hline
$1$ & $1$ & $2$ & $1$ & $2$\\\hline
$2$ & $2$ & $1$ & $0$ & $3$\\\hline
\end{tabular}\end{center}

\step{2}{Part a: largest and smallest sums}
For the largest sum, choose all $4$s in type III and all $5$s in type IV. The sum is
\[6a+3(3-a)+4(2-a)+5(1+a)=22+4a,\]
largest when $a=2$: sum $30$ (the die $3,5,5,5,6,6$).
For the smallest sum, choose all $2$s in type III and all $1$s in type IV. The sum is
\[6a+3(3-a)+2(2-a)+1\cdot(1+a)=14+2a,\]
smallest when $a=0$: sum $14$ (the die $1,2,2,3,3,3$).
\begin{recallbox}Collecting terms: $6a-3a-4a+5a=4a$ and $9+8+5=22$; similarly $6a-3a-2a+a=2a$ and $9+4+1=14$.\end{recallbox}

\step{3}{Part b: count the dice}
Once $a$ is known, we only have to decide how many of the type III faces are $4$ (the rest are $2$), and how many of the type IV faces are
$5$ (the rest are $1$).
\begin{itemize}[leftmargin=1.2em,itemsep=1pt,topsep=2pt]
\item $a=0$: type III has $2$ faces, the number of $4$s can be $0,1,2$: $3$ ways; type IV has $1$ face: $2$ ways. That gives $3\times2=6$ dice.
\item $a=1$: type III has $1$ face: $2$ ways; type IV has $2$ faces: $3$ ways. That gives $2\times3=6$ dice.
\item $a=2$: type III has $0$ faces: $1$ way; type IV has $3$ faces (the number of $5$s is $0,1,2,3$): $4$ ways. That gives $1\times4=4$ dice.
\end{itemize}
Total: $6+6+4=16$ dice.
\begin{warnbox}Since the order of the faces does not matter, two type III faces ``$2,4$'' or ``$4,2$'' are \emph{one} case. Counting the
\emph{number} of $4$s ($0$, $1$ or $2$) instead of arranging them avoids double counting.\end{warnbox}
\verify{Two dice from the list: with $a=1$, one $4$ and two $1$s: the die $1,1,3,3,4,6$ has $3$ multiples of $3$ ($3,3,6$) and $2$ even
faces ($4,6$). Correct. The die $1,2,2,3,3,3$ (sum $14$) also has $3$ multiples of $3$ and $2$ even faces.}
\ans{a) Largest $30$, smallest $14$;\quad b) $16$ dice}
}

\begin{probbox}{45}{\stars{1}}
Three lights flash: the first flashes every $4$ seconds, the second every $6$ seconds, the third every $10$ seconds. At the start
($0$ seconds) all three flash together. Consider the next $10$ minutes (from second $1$ to the end of second $600$).
\begin{enumerate}[label=\alph*),leftmargin=1.6em,itemsep=1pt,topsep=2pt]
\item How many times do all three lights flash together?
\item At how many moments do \emph{exactly two} lights flash together?
\item At how many moments does at least one light flash?
\end{enumerate}
\end{probbox}
\probsol{
\textbf{Understanding the problem.} The first light flashes at the seconds that are multiples of $4$; the second at multiples of $6$; the
third at multiples of $10$ (the same idea as in problem 20: moments that ``coincide are found with common multiples). We count multiples from $1$ to $600$.

\step{1}{Count how many times each light flashes}
\begin{recallbox}The number of multiples of $d$ among $1$ to $N$ is the whole-number part of $N\div d$ (the quotient in division with
remainder).\end{recallbox}
Light 1: $600\div4=150$ times. Light 2: $600\div6=100$ times. Light 3: $600\div10=60$ times.

\step{2}{Two and three lights together -- use the LCM}
Two lights flash together when the number of seconds is a common multiple of the two periods, that is a multiple of their LCM.
\begin{itemize}[leftmargin=1.2em,itemsep=1pt,topsep=2pt]
\item $\text{lcm}(4,6)=12$: $600\div12=50$ times lights 1 and 2 flash together.
\item $\text{lcm}(4,10)=20$: $600\div20=30$ times.
\item $\text{lcm}(6,10)=30$: $600\div30=20$ times.
\item $\text{lcm}(4,6,10)=60$: $600\div60=10$ times all three flash together. This answers part a.
\end{itemize}
\begin{warnbox}$\text{lcm}(4,6)$ is $12$, not $4\times6=24$ (because $4$ and $6$ have the common divisor $2$). Similarly
$\text{lcm}(4,6,10)=60$, not $240$.\end{warnbox}

\step{3}{Part b: exactly two lights}
Of the $50$ times lights 1 and 2 flash together, light 3 also flashes $10$ times (all three), so only lights 1 and 2 flash $50-10=40$ times.
Similarly: only lights 1 and 3: $30-10=20$; only lights 2 and 3: $20-10=10$. Total number of moments with exactly two lights:
$40+20+10=70$.

\step{4}{Part c: draw a Venn diagram and fill it in from the inside out}
Only light 1: $150-40-20-10=80$; only light 2: $100-40-10-10=40$; only light 3: $60-20-10-10=20$.
\begin{figC}
  \draw[tmBlue,thick] (0,0) circle (1.3);\draw[tmOrange,thick] (1.6,0) circle (1.3);\draw[tmGreen!80!black,thick] (0.8,-1.3) circle (1.3);
  \node[text=tmBlue,anchor=south east,inner sep=1pt] at (-0.95,0.95) {light 1};\node[text=tmOrange!80!black,anchor=south west,inner sep=1pt] at (2.55,0.95) {light 2};
  \node[text=tmGreen!70!black,anchor=north,inner sep=2pt] at (0.8,-2.62) {light 3};
  \node at (-0.45,0.25) {$80$};\node at (2.05,0.25) {$40$};\node at (0.8,-1.95) {$20$};
  \node at (0.8,0.5) {$40$};\node at (0.05,-0.95) {$20$};\node at (1.55,-0.95) {$10$};
  \node[text=AnsRed,font=\bfseries\small] at (0.8,-0.45) {$10$};
\end{figC}
The number of moments when at least one light flashes is the sum of the seven numbers in the diagram:
\[80+40+20+40+20+10+10=220.\]
\begin{recallbox}\emph{Inclusion--exclusion} for three groups: number of elements in at least one group $=$ (sum of the sizes of the three groups)
$-$ (sum of the sizes of the overlaps of each pair of groups) $+$ (size of the overlap of all three groups). The overlap of all three is added
$3$ times and subtracted $3$ times, so it must be added back once.\end{recallbox}
\verify{By inclusion--exclusion: $150+100+60-50-30-20+10=220$. It matches the Venn diagram.}
\ans{a) $10$ times;\quad b) $70$ moments;\quad c) $220$ moments}
}

\begin{probbox}{46}{\stars{2}}
\begin{enumerate}[label=\alph*),leftmargin=1.6em,itemsep=2pt,topsep=2pt]
\item A square pyramid has $5$ faces: a square base and four triangular side faces. What is the smallest number of colours needed to paint
the $5$ faces (one colour per face) so that any two faces sharing an edge have different colours?
\item There are $6$ different colours. We paint the six faces of a cube, one colour per face, using all $6$ colours. Two cubes are considered
the same if one can be rotated to look exactly like the other. How many different cubes are there?
\end{enumerate}
\begin{figC}
  \begin{scope}
    \coordinate (A) at (0,0);\coordinate (B) at (2,0);\coordinate (C) at (2.7,0.6);\coordinate (D) at (0.7,0.6);\coordinate (S) at (1.35,2.1);
    \fill[tmBlue!30] (A)--(B)--(S)--cycle;\fill[tmOrange!35] (B)--(C)--(S)--cycle;
    \draw[Ink,thick,dashed] (A)--(D)--(C);\draw[Ink,thick,dashed] (D)--(S);
    \draw[Ink,thick] (A)--(B)--(C)--(S)--cycle;\draw[Ink,thick] (B)--(S);
    \node[font=\scriptsize] at (1.35,-0.3) {square pyramid};
  \end{scope}
  \begin{scope}[shift={(4.2,0)}]
    \fill[tmGreen!35,draw=Ink,thick] (0,0)--(1.4,0)--(1.4,1.4)--(0,1.4)--cycle;
    \fill[tmPurple!30,draw=Ink,thick] (0,1.4)--(1.4,1.4)--(2.0,1.95)--(0.6,1.95)--cycle;
    \fill[tmTeal!30,draw=Ink,thick] (1.4,0)--(2.0,0.55)--(2.0,1.95)--(1.4,1.4)--cycle;
    \node[font=\scriptsize] at (1,-0.3) {cube with $6$ colours};
  \end{scope}
\end{figC}
\end{probbox}
\probsol{
\textbf{Understanding the problem.} Part a is like problem 21: find a colouring and prove that fewer colours are impossible. Part b is a
counting problem ``up to rotation'': we must be careful because a cube can be placed in many positions.

\step{1}{Part a: at least $3$ colours are needed}
\begin{proofbox}The base shares an edge with all four side faces, so its colour differs from the colour of every side face. Two neighbouring
side faces also share an edge (a slanted edge going up to the apex), so they must differ too. Take the base and two neighbouring side faces:
any two of these three faces share an edge, so at least $3$ colours are needed.\end{proofbox}
\emph{Three colours are enough:} paint the base colour $1$; going around, paint the four side faces alternately colour $2$, colour $3$,
colour $2$, colour $3$. The fourth side face (colour $3$) is next to the first (colour $2$): different. So part a needs at least $3$ colours.
\begin{tipbox}If the base were a pentagon (a pyramid with $5$ side faces), the $5$ side faces would form an odd ring and could not alternate
between two colours, so $4$ colours would be needed. Whether the number of side faces is even or odd decides the answer!\end{tipbox}

\step{2}{Part b: put the cube in a ``standard position''}
Call the colours $1,2,3,4,5,6$. Rotate each cube in two steps:
\begin{itemize}[leftmargin=1.2em,itemsep=1pt,topsep=2pt]
\item rotate it so that the face of colour $1$ is \emph{at the bottom};
\item keeping the bottom fixed, turn it about the vertical axis so that, among the four side faces, the one with the \emph{smallest colour
number} is at the front.
\end{itemize}
Once placed like this, no rotation keeps both the bottom face and the front face in place. Therefore two cubes are the same exactly when,
in standard position, they have the same colour in every position.

\step{3}{Count the standard positions}
\begin{itemize}[leftmargin=1.2em,itemsep=1pt,topsep=2pt]
\item Top face (opposite colour $1$): one of the $5$ remaining colours: $5$ ways.
\item The four remaining colours go on the four side faces; the smallest of them is already placed at the front (no choice).
\item The other three colours go on the three faces left, back, right: $3\times2\times1=6$ ways.
\end{itemize}
By the multiplication principle: $5\times6=30$ different cubes.

\step{4}{Check with a second method}
If rotations were \emph{not} allowed (the six positions top, bottom, front, back, left, right fixed), the number of paintings would be
$6\times5\times4\times3\times2\times1=720$. Each cube can be placed in $6\times4=24$ positions ($6$ choices for the face at the bottom, then
$4$ turns about the vertical axis), and with $6$ different colours these $24$ positions give $24$ different fixed paintings.
So the number of different cubes is $720\div24=30$, matching the count in step 3.
\begin{warnbox}Answering $720$ forgets that cubes which can be rotated into each other are the same. Answering $6\times5=30$ ``by luck''
without explanation is not enough either: you must argue why each cube is counted exactly once.\end{warnbox}
\ans{a) $3$ colours;\quad b) $30$ cubes}
}

\begin{probbox}{47}{\stars{2}}
Consider the integers $1,2,3,\dots,20$.
\begin{enumerate}[label=\alph*),leftmargin=1.6em,itemsep=1pt,topsep=2pt]
\item How many pairs $(m,n)$ with $m<n$ are there such that $m+n$ is divisible by $3$?
\item How many pairs $(m,n)$ with $m<n$ are there such that $m+n$ is divisible by $4$?
\item How many triples $(m,n,p)$ with $m<n<p$ are there such that $m+n+p$ is divisible by $3$?
\end{enumerate}
\end{probbox}
\probsol{
\textbf{Understanding the problem.} In problem 22 we sorted the numbers by parity (the remainder on division by $2$). Here we sort them by
\emph{the remainder on division by $3$} (or by $4$). Because $m<n$ (and $m<n<p$), each pair or triple is counted once, regardless of order.
\begin{recallbox}The remainder of a sum on division by $3$ equals the remainder of the sum of the remainders. For example $7+11=18$: $7$ leaves
remainder $1$, $11$ leaves remainder $2$, $1+2=3$ is divisible by $3$, so $18$ is divisible by $3$.\end{recallbox}

\step{1}{Sort the $20$ numbers by remainder on division by $3$}
\begin{center}\begin{tabular}{|c|l|c|}\hline
Remainder & Numbers & How many\\\hline
$0$ & $3,6,9,12,15,18$ & $6$\\\hline
$1$ & $1,4,7,10,13,16,19$ & $7$\\\hline
$2$ & $2,5,8,11,14,17,20$ & $7$\\\hline
\end{tabular}\end{center}

\step{2}{Part a}
The sum of two remainders is divisible by $3$ in only two cases: $0+0$ or $1+2$.
\begin{itemize}[leftmargin=1.2em,itemsep=1pt,topsep=2pt]
\item Both numbers with remainder $0$: choose $2$ of the $6$ numbers: $\dfrac{6\times5}{2}=15$ pairs.
\item One with remainder $1$, one with remainder $2$: $7\times7=49$ pairs (the numbers are in different groups, so automatically different;
the smaller one is $m$).
\end{itemize}
(Two numbers with remainder $1$ give a sum with remainder $2$; two numbers with remainder $2$ give $4$, remainder $1$: not allowed.)
Total: $15+49=64$ pairs.

\step{3}{Part b: sort by remainder on division by $4$}
Each remainder $0,1,2,3$ has exactly $5$ numbers (because $20=4\times5$). The sum of two remainders is divisible by $4$ for: $0+0$, $2+2$
or $1+3$.
\[\frac{5\times4}{2}+\frac{5\times4}{2}+5\times5=10+10+25=45 \text{ pairs.}\]

\step{4}{Part c: the remainder patterns of a triple}
List every pattern of remainders (order not important) of three numbers and its sum:
\begin{center}\begin{tabular}{|c|c|c|c|c|c|c|c|c|c|c|}\hline
Pattern & $000$ & $111$ & $222$ & $012$ & $001$ & $002$ & $011$ & $022$ & $112$ & $122$\\\hline
Sum & $0$ & $3$ & $6$ & $3$ & $1$ & $2$ & $2$ & $4$ & $4$ & $5$\\\hline
Divisible by $3$? & yes & yes & yes & yes & no & no & no & no & no & no\\\hline
\end{tabular}\end{center}
Count the four good patterns:
\begin{itemize}[leftmargin=1.2em,itemsep=1pt,topsep=2pt]
\item Pattern $000$: choose $3$ of $6$ numbers: $\dfrac{6\times5\times4}{3\times2\times1}=20$.
\item Pattern $111$: choose $3$ of $7$ numbers: $\dfrac{7\times6\times5}{3\times2\times1}=35$. Pattern $222$: also $35$.
\item Pattern $012$: one number from each group: $6\times7\times7=294$.
\end{itemize}
Total: $20+35+35+294=384$ triples.
\begin{recallbox}The number of ways to choose $3$ objects (order not important) from $k$ objects is $\dfrac{k\times(k-1)\times(k-2)}{6}$:
there are $k(k-1)(k-2)$ ordered choices, and each triple is counted $3\times2\times1=6$ times.\end{recallbox}
\begin{warnbox}For pattern $012$ we do \emph{not} divide by $6$: the three numbers come from three different groups, so each choice (one number
with remainder $0$, one with remainder $1$, one with remainder $2$) gives exactly one triple, which we then write in increasing order.\end{warnbox}
\verify{The total number of triples is $\dfrac{20\times19\times18}{6}=1140$. There are $384$ triples with sum divisible by $3$, close to
$\dfrac13$ of $1140$ ($380$), which is reasonable since the sums are spread ``fairly evenly'' among the three remainders.}
\ans{a) $64$ pairs;\quad b) $45$ pairs;\quad c) $384$ triples}
}

\begin{probbox}{48}{\stars{1}}
For a two-digit number $n=\nd{ab}$, call $\nd{ba}$ the \emph{reversal} of $n$ (as in problem 23).
\begin{enumerate}[label=\alph*),leftmargin=1.6em,itemsep=1pt,topsep=2pt]
\item Find all two-digit numbers $n$ such that $n$ plus its reversal is a perfect square.
\item Find all two-digit numbers with both digits non-zero such that $n$ minus its reversal is a positive perfect square.
\item For a three-digit number $\nd{abc}$ ($a\ne0$, $c\ne0$), prove that the difference $\nd{abc}-\nd{cba}$ is always divisible by $99$.
\end{enumerate}
\end{probbox}
\probsol{
\textbf{Understanding the problem.} Write the numbers in terms of their digits: $\nd{ab}=10a+b$ and $\nd{ba}=10b+a$, where $a$ is the tens digit
($1\le a\le9$) and $b$ is the units digit. A perfect square is the square of a natural number: $1,4,9,16,25,\dots$

\step{1}{Part a: compute the sum}
$\nd{ab}+\nd{ba}=(10a+b)+(10b+a)=11a+11b=11(a+b)$.
\begin{proofbox}If $11(a+b)=s\times s$, then $s\times s$ is divisible by the prime $11$, so $s$ is divisible by $11$ (if a product of two numbers
is divisible by a prime $p$, one of the two numbers is divisible by $p$). Then $s\times s$ is divisible by $121$, so $11(a+b)$ is divisible by
$121$, and $a+b$ is divisible by $11$.\end{proofbox}
For a two-digit number (the reversal may start with $0$, e.g.\ $\nd{ba}=b$ when $b=0$) we have $1\le a+b\le18$, so $a+b=11$, and then the
sum is $11\times11=121=11^2$, which is indeed a perfect square. The numbers are those with $a+b=11$, $a$ from $2$ to $9$:
\[29,\ 38,\ 47,\ 56,\ 65,\ 74,\ 83,\ 92 \quad(8 \text{ numbers}).\]

\step{2}{Part b: compute the difference}
$\nd{ab}-\nd{ba}=(10a+b)-(10b+a)=9a-9b=9(a-b)$. The difference is positive, so $a>b\ge1$.
Since $9=3\times3$ is a perfect square, $9(a-b)$ is a perfect square exactly when $a-b$ is a perfect square.
\begin{proofbox}If $a-b=t\times t$, then $9(a-b)=(3t)\times(3t)$. Conversely, if $9(a-b)=s\times s$, then $s\times s$ is divisible by $3$, so
$s$ is divisible by $3$; writing $s=3t$ gives $9(a-b)=9t\times t$, that is $a-b=t\times t$.\end{proofbox}
Since $1\le b<a\le9$, we have $1\le a-b\le8$; the perfect squares in this range are $1$ and $4$.
\begin{itemize}[leftmargin=1.2em,itemsep=1pt,topsep=2pt]
\item $a-b=1$: $21, 32, 43, 54, 65, 76, 87, 98$ ($8$ numbers).
\item $a-b=4$: $51, 62, 73, 84, 95$ ($5$ numbers).
\end{itemize}
There are $13$ numbers in all.

\step{3}{Part c: split by place value}
Write both numbers using place values: $\nd{abc}=100a+10b+c$ and $\nd{cba}=100c+10b+a$. Then
\[\nd{abc}-\nd{cba}=(100a+10b+c)-(100c+10b+a)=100a-a+10b-10b+c-100c,\]
\[\text{that is}\quad\nd{abc}-\nd{cba}=99a-99c=99(a-c).\]
Since $a-c$ is an integer, $99(a-c)$ is divisible by $99$. (The tens digit $b$ cancels out completely!)
\begin{warnbox}$\nd{ab}$ denotes the number with digits $a$ and $b$, \emph{not} the product $a\times b$. Always rewrite $\nd{ab}$ as $10a+b$
before calculating.\end{warnbox}
\verify{a) $47+74=121=11^2$. b) $73-37=36=6^2$, $98-89=9=3^2$. c) $731-137=594=99\times6$; $205-502=-297=99\times(-3)$.}
\ans{a) $29$, $38$, $47$, $56$, $65$, $74$, $83$, $92$;\quad b) $13$ numbers (listed in step 2);\quad \mbox{c) $\nd{abc}-\nd{cba}=99(a-c)$.}}
}

\needspace{16\baselineskip}
\begin{probbox}{49}{\stars{2}}
A grid of $3$ rows and $5$ columns of squares ($15$ small squares) is shown; the square in the very centre is shaded.
\begin{figC}\begin{scope}[scale=0.6]
  \fill[tmOrange!45] (2,1) rectangle (3,2);
  \draw[Ink,thick] (0,0) grid (5,3);
\end{scope}\end{figC}
\begin{enumerate}[label=\alph*),leftmargin=1.6em,itemsep=1pt,topsep=2pt]
\item How many rectangles (including squares) with sides on the grid lines are there in the figure?
\item How many of them are squares?
\item One of the rectangles in part a is chosen at random. Find the probability that it contains the shaded square.
\end{enumerate}
\end{probbox}
\probsol{
\textbf{Understanding the problem.} As in problem 24, each rectangle is determined by two horizontal grid lines (top and bottom sides) and
two vertical grid lines (left and right sides). A grid with $3$ rows has $4$ horizontal lines, numbered $0,1,2,3$ from top to bottom; a grid
with $5$ columns has $6$ vertical lines, numbered $0,1,\dots,5$ from left to right.

\step{1}{Part a: count the rectangles}
Choose $2$ of the $4$ horizontal lines: $\dfrac{4\times3}{2}=6$ ways. Choose $2$ of the $6$ vertical lines: $\dfrac{6\times5}{2}=15$ ways.
By the multiplication principle: $6\times15=90$ rectangles.

\step{2}{Part b: count the squares by size}
\begin{recallbox}In a $3\times5$ grid, a square of side $k$ has $(3-k+1)$ positions vertically and $(5-k+1)$ positions horizontally.\end{recallbox}
\begin{itemize}[leftmargin=1.2em,itemsep=1pt,topsep=2pt]
\item side $1$: $3\times5=15$ squares;
\item side $2$: $2\times4=8$ squares;
\item side $3$: $1\times3=3$ squares;
\item side $4$ or more: none (the grid is only $3$ high).
\end{itemize}
Total: $15+8+3=26$ squares.
\begin{warnbox}Do not use ``choose $2$ horizontal and $2$ vertical lines'' to count squares: most of those choices give rectangles that are not
squares. Squares must be counted by side length.\end{warnbox}

\step{3}{Part c: rectangles containing the centre square}
The centre square lies between horizontal lines $1$ and $2$, and between vertical lines $2$ and $3$. A rectangle contains it exactly when:
\begin{itemize}[leftmargin=1.2em,itemsep=1pt,topsep=2pt]
\item its top side is horizontal line $0$ or $1$ (above the square): $2$ ways; its bottom side is horizontal line $2$ or $3$: $2$ ways;
\item its left side is vertical line $0$, $1$ or $2$: $3$ ways; its right side is vertical line $3$, $4$ or $5$: $3$ ways.
\end{itemize}
\begin{figC}\begin{scope}[scale=0.6]
  \fill[tmGreen!25] (1,0) rectangle (4,2);
  \fill[tmOrange!55] (2,1) rectangle (3,2);
  \draw[Ink,thick] (0,0) grid (5,3);
  \draw[tmGreen!70!black,line width=2pt] (1,0) rectangle (4,2);
  \foreach \i in {0,...,5}{\node[font=\scriptsize,text=Slate] at (\i,3.35) {\i};}
  \foreach \j/\l in {3/0,2/1,1/2,0/3}{\node[font=\scriptsize,text=Slate] at (-0.35,\j) {\l};}
\end{scope}\end{figC}
(For example, the green rectangle: top side horizontal line $1$, bottom side horizontal line $3$, left side vertical line $1$, right side
vertical line $4$.)
Number of rectangles containing the centre square: $2\times2\times3\times3=36$.
\[P=\frac{36}{90}=\frac{2}{5}.\]
\begin{proofbox}Why is this condition right? A rectangle contains the centre square exactly when it ``covers'' that square in both directions:
its top side must be above (or on) the top side of the square, its bottom side below (or on) the bottom side of the square; similarly for the
left and right sides. These four choices are independent, so we multiply.\end{proofbox}
\verify{Split by direction: of the $15$ choices of two vertical lines, $3\times3=9$ ``cover'' the middle column; of the $6$ choices of two
horizontal lines, $2\times2=4$ cover the middle row. Probability $=\dfrac{9}{15}\times\dfrac{4}{6}=\dfrac35\times\dfrac23=\dfrac25$. It matches.}
\ans{a) $90$;\quad b) $26$;\quad c) $\dfrac25$}
}

\begin{probbox}{50}{\stars{2}}
A box contains $2026$ red marbles and $2027$ blue marbles. Outside the box there is a very large pile of blue marbles. Felix repeats the
following step: take $2$ marbles from the box at random.
\begin{itemize}[leftmargin=1.2em,itemsep=1pt,topsep=2pt]
\item If the two marbles have the \emph{same colour} (both red or both blue): throw both away, then take $1$ blue marble from the pile
outside and put it into the box.
\item If the two marbles have \emph{different colours}: throw away the blue one and put the red one back into the box.
\end{itemize}
Felix stops when only $1$ marble is left in the box.
\begin{enumerate}[label=\alph*),leftmargin=1.6em,itemsep=1pt,topsep=2pt]
\item How many steps did Felix carry out?
\item Find the probability that the last marble is red.
\item If the box had started with $2025$ red marbles and $2027$ blue marbles, what would the probability in part b be?
\end{enumerate}
\end{probbox}
\probsol{
\textbf{Understanding the problem.} There are thousands of marbles, so we cannot draw a probability tree as in problem 25. Instead we look for
a quantity that \emph{does not change} in any step (called an \emph{invariant}). In problem 25 the invariant was ``there is always a blue
marble''; here we look at \emph{the parity of the number of red marbles}.

\step{1}{Part a: each step reduces the number of marbles by exactly $1$}
\begin{itemize}[leftmargin=1.2em,itemsep=1pt,topsep=2pt]
\item Same colour: $2$ taken out, $1$ put in $\Rightarrow$ the number of marbles in the box decreases by $1$.
\item Different colours: $2$ taken out, $1$ put back $\Rightarrow$ the number decreases by $1$.
\end{itemize}
At the start there are $2026+2027=4053$ marbles, and at the end $1$ marble is left. The number of steps is $4053-1=4052$.

\step{2}{Follow the number of red marbles in each case}
\begin{center}\begin{tabular}{|c|c|c|}\hline
Drawn & Change in red marbles & Change in blue marbles\\\hline
red, red & decreases by $2$ & increases by $1$\\\hline
blue, blue & no change & down $2$, up $1$: decreases by $1$\\\hline
red, blue & no change (the red is put back) & decreases by $1$\\\hline
\end{tabular}\end{center}
\begin{proofbox}The number of red marbles either stays the same or decreases by exactly $2$. Subtracting $2$ does not change parity (an even
number minus $2$ is still even, an odd number minus $2$ is still odd). So \textbf{the parity of the number of red marbles never changes}, no
matter which marbles Felix draws.\end{proofbox}

\step{3}{Part b: $2026$ red marbles at the start (an even number)}
The number of red marbles is always even. When only $1$ marble is left: if it were red, the number of red marbles would be $1$ (odd) --
impossible. So the last marble is certainly blue. The event ``the last marble is red'' is an impossible event:
\[P(\text{last marble red})=0.\]

\step{4}{Part c: $2025$ red marbles at the start (an odd number)}
The number of red marbles is always odd, so it is never $0$. When $1$ marble is left, the number of red marbles is an odd number not exceeding
$1$, that is $1$: the last marble is certainly red. This is a certain event:
\[P(\text{last marble red})=1.\]
\begin{tipbox}The initial number of blue marbles has no effect on the result! Only the parity of the number of red marbles decides the colour
of the last marble.\end{tipbox}
\begin{parentbox}Play with your child using a few marbles (for example $2$ red and $3$ blue, or $3$ red and $2$ blue), or with coins of two
colours. Seeing the same result every time makes children curious about ``why''; then give a hint: ``Try counting the red marbles after each
turn: is the number even or odd?''\end{parentbox}
\begin{warnbox}Seeing the word ``random'' and thinking the answer must be a ``nice'' fraction such as $\frac12$ is a mistake. Some random
processes have a final result that is completely determined.\end{warnbox}
\verify{Try $2$ red, $1$ blue: drawing red--red $\to(0\text{ red},2\text{ blue})\to(0,1)$: blue; drawing red--blue $\to(2,0)\to$ draw
red--red $\to(0,1)$: blue. Every way ends with a blue marble, consistent with an even initial number of red marbles.}
\ans{a) $4052$ steps;\quad b) $0$;\quad c) $1$}
}
